\scanpage{6}
\mspage{5}
Alors $G'=P\times^G G=\operatorname{ad}(P)$

En fait, il y a des cas intéressants où \deleted{si} on peut trouver aucun foncteur fibre. Mais si $C$ est de type fini, alors on \inserted{doit pouvoir} montrer qu'on peut trouver un foncteur fibre sur une extension finie de $K$. \deleted{On \unsure{réussirait|construit} ainsi à décrire les} \inserted{Quand il en est ainsi, \deleted{on \unsure{doit avoir|démontre}}} admettant le résultat ``d'unicité'' de foncteurs fibres, \deleted{on doit} arriver à décrire $(C,F)$ intrinsèquement en termes d'une gerbe sur un site convenable\ldots

8. \uline{Cas des catégories des motifs}. \deleted{\uline{Lien avec le groupe fondamental profini}.}

$S$ schéma régulier \inserted{\unsure{rédiger|intègre|connexe}}. On définit $M^+(S)$, $M(S)$ en paraphrasant ce qui est fait sur un corps, en travaillant (pour $X$ projectif lisse sur $S$, avec $C(X)\overset{\mathrm{dfn}}=C(X_\eta)$ ($\eta$ pt générique de $S$). Ainsi, $M(S)\unsure{\longrightarrow|-}M(\eta)$ est pleinement fidèle, permettant d'interpréter $M(S)$ comme ss-catégorie pleine de $M(\eta)$. De plus, pour tt morphisme $f:S'\longrightarrow S$ de schémas réguliers, on déduit $f^*:M(S)\longrightarrow M(S')$ un $\otimes$-foncteur correspondant (c'est ici que sert l'hypothèse de \uline{régularité} sur les schémas\ldots.).
% slanted
\margin{à expliciter\ldots.}

Tout point \inserted{géom} \uline{complexe} $\xi$ de $S$ définit un foncteur fibre ``de Betti'' \deleted{$\unsure{F^\xi_{\mathrm{Betti}}=F_\xi}$} $F^{\mathrm{Betti}}_\xi=F^B_\xi:M(S)\unsure{-|\longrightarrow}\mathrm{Modf}(\unsure{Q|\mathbb Q})$, d'où \unsure{un|une} groupe proalgébrique $G$ sur $Q$, proréductif
% archive sheet number: 5

